\section{Introduction}

\subsection{Scope and Objectives}

This take-home report studies five natural language understanding and generation problems that cover both analysis and generation settings: sentiment classification, named entity recognition, summarization, machine translation, and language modeling. The common objective across all questions is not only to obtain strong final metrics, but also to compare different representation regimes in a controlled way. Wherever possible, the experiments contrast classical feature-based baselines, supervised neural architectures, and transformer or pretrained contextual models so that the final discussion can separate the effect of model class from the effect of representation quality.

This report presents a complete five-task experimental study built on verified run artifacts. Questions 1 and 2 provide full multi-model comparisons over completed train, validation, and test exports; Question 3 reports a controlled classical-versus-pretrained summarization comparison on a fixed evaluation slice; and Questions 4 and 5 now report full-data neural runs with saved training histories and qualitative artifacts. This organization keeps all claims tied to reproducible outputs and avoids speculative results.

\subsection{Technical Setup}

All experiments are managed from YAML configuration files, use a fixed random seed of 42 for reproducibility, and write timestamped run directories that include the resolved configuration, package environment snapshot, metrics, and analysis artifacts. The implementation stack combines Hugging Face Datasets for dataset access, scikit-learn and related linear modeling tools for classical baselines, PyTorch for custom neural models, and transformers for pretrained contextual encoders. Report assembly follows the same principle of artifact stability: figures copied into \texttt{report/figures/} and tables placed in \texttt{report/tables/} are sourced from completed experiment outputs rather than ad hoc terminal observations.

Table~\ref{tab:introduction_task_overview} summarizes the five questions, their datasets, and the representative model families covered by the report.

\begin{table}[H]
\centering
\footnotesize
\begin{tabularx}{\linewidth}{c l l >{\raggedright\arraybackslash}X}
\toprule
Q & Task & Dataset & Representative Models \\
\midrule
1 & Text classification & IMDb & TF-IDF+LR/SVM, BiLSTM, DistilBERT \\
2 & Named entity recognition & CoNLL-2003 & CRF, BiLSTM-CRF, BERT \\
3 & Summarization & CNN/DailyMail & TextRank, DistilBART \\
4 & Machine translation & Multi30k & Seq2Seq+Attn, Opus-MT \\
5 & Language modeling & WikiText-2 & Trigram, LSTM, distilGPT-2 \\
\bottomrule
\end{tabularx}
\caption{Task and dataset overview for the take-home report.}
\label{tab:introduction_task_overview}
\end{table}

\subsection{Report Organization}

Section 2 presents Question 1, which analyzes representation learning for IMDb sentiment classification and compares TF-IDF baselines, a word-level BiLSTM, and DistilBERT on the full exported split. Section 3 presents Question 2, which studies named entity recognition on CoNLL-2003 with CRF, BiLSTM-CRF, and BERT token classification. Section 4 presents Question 3, which reports a direct TextRank-versus-DistilBART summarization comparison on a fixed exported evaluation slice. Section 5 presents Question 4, which reports a full-data seq2seq-plus-attention translation run on Multi30k together with attention visualizations. Section 6 presents Question 5, which reports a full-data LSTM language-model run on WikiText-2 with training curves and generation samples. The conclusion section synthesizes cross-task findings and methodological lessons for fair NLP model comparison.